%% kapa: exam=1 round=20 year=2009 period=2 form=A subject=appraisal_law subject_ko=부동산관계법규 seq=35 number=75
%% source: Q-Net 감정평가사 기출문제 (cst00309) artlSeq=5024178 file=extracted/1008/5024178/2009년 제20회 감정평가사 1차시험 기출문제/2교시.hwp
%% conversion: kordoc 4.24.51 -> tools/gen1.py (auto); review=proofread-sonnet; flags=none
\begin{question}{75}
지적법령상 지적측량에 의하여 경계점 및 좌표를 결정하는 경우가 \uline{아닌} 것은?
\begin{choices}
  \item 신규등록을 위한 경계점 및 좌표의 결정
  \item 등록전환을 위한 경계점 및 좌표의 결정
  \item 분할을 위한 경계점 및 좌표의 결정
  \item 합병을 위한 경계점 및 좌표의 결정
  \item 축척변경을 위한 경계점 및 좌표의 결정
\end{choices}
\end{question}
\answer{4}
